\documentclass[10pt,a4paper,withhyper]{altacv}

\usepackage{paracol}

\iftutex
  \setmainfont{Roboto Slab}
  \setsansfont{Lato}
  \renewcommand{\familydefault}{\sfdefault}
\else
  \usepackage[rm]{roboto}
  \usepackage[defaultsans]{lato}
  \renewcommand{\familydefault}{\sfdefault}
\fi

% ===== CORES =====
\definecolor{DarkBlue}{HTML}{0B1F3A}
\definecolor{MidBlue}{HTML}{163A6B}
\definecolor{LightBlue}{HTML}{4F81BD}
\definecolor{BodyGrey}{HTML}{444444}

\colorlet{name}{black}
\colorlet{tagline}{MidBlue}
\colorlet{heading}{DarkBlue}
\colorlet{headingrule}{LightBlue}
\colorlet{subheading}{MidBlue}
\colorlet{accent}{LightBlue}
\colorlet{body}{BodyGrey}

\renewcommand{\namefont}{\Huge\bfseries}
\renewcommand{\cvsectionfont}{\Large\bfseries}

\begin{document}

\name{João Pedro Miranda Câmara}
\tagline{Graduando em Astronomia — UFRJ}

\personalinfo{
  \email{joao24@ov.ufrj.br}
  \location{Rio de Janeiro, Brasil}
  \xtwitter{@Pedrottsx}
  \printinfo{\faInstagram}{\href{https://instagram.com/jpcamaracrvg}{@jpcamaracrvg}}
  \linkedin{jpcamara}
}

% SATURNO
\photoR{2.8cm}{saturno.jpg}

\makecvheader

\cvsection{Projetos}

\cvevent{Criação de um eletroscópio caseiro}{}{}{}

\begin{itemize}
\item Experimento demonstrando eletrização por atrito e indução eletrostática.
\end{itemize}

\divider

\cvevent{Apresentação sobre Física Quântica}{}{}{}

Apresentação realizada durante o Ensino Médio abordando conceitos introdutórios de Física Quântica.

\cvsection{Habilidades}

\cvtag{Python}
\cvtag{Cálculo}
\cvtag{Criatividade}

\medskip

\cvsection{Idiomas}

\cvskill{Português}{5}
\divider
\cvskill{Inglês}{4}
\divider
\cvskill{Espanhol}{2.5}

\cvsection{Formação}

\cvevent
{Graduando em Astronomia}
{UFRJ}
{Jan 2024 -- Atual}
{Rio de Janeiro}
\cvsection{Referências}

Prof.\ Bruno Morgado — Observatório do Valongo (UFRJ), Rio de Janeiro, Brasil

\end{document}